\documentclass{article}
\usepackage[utf8]{inputenc}

\title{tarea 3}
\author{iabin Arteaga}
\date{October 2017}

\usepackage{natbib}
\usepackage{graphicx}

\begin{document}

\maketitle

\section{4.1}
Vamos a suponer que las funciones, suma y producto de naturales son equivalentes en ambos lenguajes
\\Por Demostrar que ambas expresiones son equivalentes\\

Lo demostraremos por inducción.\\\\
Caso base
\\$n = 0$
\\Cuando ejecutamos f(0)=1 y g(0)=1 , en ambas funciones es trivial ver que $1 = 1$ y son iguales\\
\\
Caso inductivo\\
Suponemos que para n , las expresiones son equivalentes, osea que ouput() de ambos programas con n son equivalentes\\\\
Por Demostrar  f(n+1) es equivalente a g(n+1)\\
\\La salida de ambos programas es (n+1)*factorial(n), pero las operaciones suma y producto son equivalentes en ambos lenguajes
además, por hipótesis de inducción, f(n)  y g(n) son equivalentes, por lo tanto,  $(n+1)*factorial(n)= (n+1)*factorial(n)$ en ambos lenguajes
\\YA

\section{4.2}
Sean L1 y L2 lenguajes Turing-completos\\
Sea E1 una expresión sobre L1\\
Sea P1 programa que contenga esta expresión. \\\\

Por definición, existe una función f que recibe un programa P sobre un lenguaje Turing-completo y devuelve una máquina de turing M equivalente
\\Pero como L2 es turing completo, puede modelar cualquier máquina de turing, en especial la máquina M , \\sea L2 de evaluación perezosa.
Entonces existe un programa P2 en L2 tal que modela la máquina M\\\\
YA

\section{4.3}
Sea F la función principal tal que el programa P1 es equivalente a F\\
Sean h1,h2...,hn funciones diferentes a F, h1(i)= a, h2(j) = b...,hn(k) = z tal que los parametros de F sean $hi$  $ \forall i <n+1$ donde n es el número de parametros de F\\
\\La diferencia entre evaluaciones es que la función F
cargaría hi(j) sería evaluación perezosa,mientras que la evaluación glotona cargaría a 'a' en el cuerpo,
pero SPDG $hi(j) = a $\\
Pero como hi es parámetro de F, por lo tanto tiene que evaluarse, así que no afecta en la evaluación de la función principal\\
Por lo tanto no afecta.\\ \\
YA

\section{Turing completez}
Debemos de probar que nuestro lenguaje puede modelar las  funciones $μ-recursivas$, funciones de los naturales a los naturales\\
1.-función constante $f(x1,x2...,xn) = k$ para toda n y k naturales
\\2.-función sucesor, creo que es obvia.
\\3.-Función de proyección, para  $P(x1,x2...,xk)=xi$  $0<i<k+1$
\\4.- Operador de sustición, es el with claramente\\
5.-Operador de recursión, el cual es obvio que tiene FWAE
\\6.-Operador de minimización, esta es la que puede causar problemas, porque necesitas poder definir condicionales,
para que puedas definir funciones que regresen 0 y generar casos base, para asegurar de que la función termine.\\
Por lo visto en clase, se vió que FWAE tiene condicionales, con esto es el último punto, y FWAE es turing completo

\end{document}
